\documentclass[materials,article,submit,moreauthors,pdftex]{Definitions/mdpi}

% MDPI internal commands
\firstpage{1}
\makeatletter
\setcounter{page}{\@firstpage}
\makeatother
\pubvolume{1}
\issuenum{1}
\articlenumber{0}
\pubyear{2026}
\copyrightyear{2026}
\datereceived{}
\daterevised{}
\dateaccepted{}
\datepublished{}
\hreflink{https://doi.org/}

\usepackage{graphicx}
\usepackage{amsmath}
\usepackage{amssymb}
\usepackage{booktabs}
\usepackage{multirow}
\usepackage{siunitx}
\usepackage{xcolor}

%=================================================================
% Title
%=================================================================
\Title{HIF-1$\alpha$-Mediated Mechanotransduction in Aortic Smooth Muscle Cells: A Multi-Layer PDMS Microfluidic Platform for High-Throughput Biomechanical Analysis}

\TitleCitation{HIF-1$\alpha$-Mediated Mechanotransduction in Aortic SMCs}

%=================================================================
% Authors
%=================================================================
\Author{Firstname Lastname $^{1,*}$, Firstname Lastname $^{2}$ and Firstname Lastname $^{1,2}$}

\AuthorNames{Firstname Lastname, Firstname Lastname and Firstname Lastname}

\AuthorCitation{Lastname, F.; Lastname, F.; Lastname, F.}

%=================================================================
% Affiliations
%=================================================================
\address{%
$^{1}$ \quad Department of Biomedical Engineering, University Name, City, Country; email@university.edu \\
$^{2}$ \quad Department of Cardiovascular Medicine, Hospital Name, City, Country}

\corres{Correspondence: corresponding.author@university.edu; Tel.: +XX-XXX-XXX-XXXX}

%=================================================================
% Abstract
%=================================================================
\abstract{Vascular smooth muscle cells (VSMCs) play a critical role in maintaining arterial homeostasis, and their dysfunction is closely associated with cardiovascular diseases including aortic calcification and aneurysm formation. Hypoxia-inducible factor 1-alpha (HIF-1$\alpha$) has emerged as a key regulator of VSMC phenotypic switching and mechanotransduction under pathological conditions. However, the biomechanical behavior of VSMCs under controlled hypoxic and mechanical stimuli remains poorly understood due to limitations in conventional culture systems. In this study, we developed a novel multi-layer polydimethylsiloxane (PDMS) microfluidic platform integrating pneumatic valve control, real-time observation capabilities, and vacuum-assisted mechanical stretching for high-throughput analysis of aortic smooth muscle cell (SMC) biomechanics. We characterized HIF-1$\alpha$ expression in human aortic tissue and blood samples from patients with aortic diseases using immunohistochemistry, Western blotting, and quantitative proteomics. Our results demonstrate that HIF-1$\alpha$ upregulation correlates with increased expression of osteogenic markers (BSP, OCN) and altered VSMC contractile phenotype. The PDMS microfluidic chip, featuring a four-layer architecture with observation channels, medium perfusion, and vacuum-driven cyclic stretch, enables simultaneous culture and biomechanical assessment of patient-derived SMCs. This platform provides a physiologically relevant model for studying the mechanical behavior of biopolymer-based vascular constructs and offers potential for personalized therapeutic screening in cardiovascular disease.}

\keyword{HIF-1$\alpha$; vascular smooth muscle cells; microfluidic; PDMS; mechanotransduction; biopolymer; aortic calcification; organ-on-chip}

%=================================================================
\begin{document}

%=================================================================
% 1. Introduction
%=================================================================
\section{Introduction}

Cardiovascular diseases (CVDs) remain the leading cause of mortality worldwide, with aortic pathologies including atherosclerosis, calcification, and aneurysm representing significant clinical challenges \cite{roth2020global}. Vascular smooth muscle cells (VSMCs) constitute the predominant cellular component of the arterial media and play essential roles in maintaining vascular tone, structural integrity, and mechanical compliance \cite{owens2004molecular}. Under pathological conditions, VSMCs undergo phenotypic modulation from a contractile to a synthetic state, characterized by reduced expression of contractile markers ($\alpha$-smooth muscle actin, SM22$\alpha$) and acquisition of proliferative, migratory, and osteogenic capabilities \cite{bennett2016vascular}.

Hypoxia-inducible factor 1-alpha (HIF-1$\alpha$) is a master transcriptional regulator that orchestrates cellular responses to hypoxic stress \cite{semenza2012hypoxia}. In VSMCs, HIF-1$\alpha$ activation has been implicated in multiple pathological processes, including vascular calcification, where it promotes osteogenic differentiation through upregulation of bone-related proteins such as bone sialoprotein (BSP) and osteocalcin (OCN) \cite{mokas2016hypoxia}. Recent studies have demonstrated that HIF-1$\alpha$ accumulation in VSMCs occurs not only under hypoxic conditions but also in response to mechanical stress, inflammatory cytokines, and metabolic dysregulation \cite{gorlach2001thrombin}. Understanding the mechanistic links between HIF-1$\alpha$ signaling and VSMC mechanical behavior is therefore crucial for developing targeted therapeutic interventions.

The mechanical properties of the vascular wall are determined by the complex interplay between cellular components and extracellular matrix (ECM) biopolymers, including collagen, elastin, and proteoglycans \cite{humphrey2014mechanotransduction}. VSMCs respond to mechanical cues through mechanotransduction pathways that convert physical stimuli into biochemical signals, influencing cell phenotype, proliferation, and matrix production \cite{haga2007molecular}. However, studying these biomechanical interactions in conventional two-dimensional culture systems fails to recapitulate the three-dimensional architecture and dynamic mechanical environment of native blood vessels.

Microfluidic organ-on-chip technology has emerged as a powerful platform for modeling vascular physiology and pathology under controlled biomechanical conditions \cite{bhatia2014microfluidic}. Polydimethylsiloxane (PDMS), a biocompatible silicone elastomer, is the most widely used material for fabricating microfluidic devices due to its optical transparency, gas permeability, tunable mechanical properties, and compatibility with soft lithography techniques \cite{mcdonald2002poly}. The mechanical behavior of PDMS as a biopolymer material directly influences cell-substrate interactions and can be engineered to mimic the stiffness of native vascular tissue \cite{palchesko2012development}.

Several microfluidic vascular models incorporating VSMCs have been developed, including systems that enable co-culture with endothelial cells under hemodynamic shear stress \cite{van2018biomimetic}, cyclic mechanical stretch \cite{zheng2012microfluidic}, and controlled oxygen tension \cite{huh2010reconstituting}. However, most existing platforms lack the capability for high-throughput analysis and integration of multiple mechanical stimuli simultaneously. Furthermore, limited studies have addressed the specific role of HIF-1$\alpha$-mediated mechanotransduction in patient-derived VSMCs using microfluidic technology.

In this study, we present a comprehensive investigation of HIF-1$\alpha$ expression and function in human aortic tissue and SMCs, combined with the development of a novel multi-layer PDMS microfluidic platform for high-throughput biomechanical analysis. Our approach integrates: (1) clinical sample analysis using immunohistochemistry, Western blotting, and quantitative proteomics to characterize HIF-1$\alpha$ pathway activation in aortic disease; (2) isolation and culture of primary human aortic SMCs; and (3) design and fabrication of a microfluidic chip with pneumatic valve control, observation channels, medium perfusion, and vacuum-driven mechanical stretching capabilities. This platform enables simultaneous exposure of SMCs to multiple physiologically relevant stimuli while permitting real-time monitoring of cellular responses, offering new insights into the mechanical behavior of biopolymer-based vascular constructs under controlled conditions.

%=================================================================
% 2. Materials and Methods
%=================================================================
\section{Materials and Methods}

\subsection{Patient Sample Collection and Ethics}

Human aortic tissue specimens were obtained from patients undergoing cardiac surgery at [Hospital Name] between [Date Range]. Peripheral blood samples were collected from the same patients preoperatively. The study protocol was approved by the Institutional Review Board (IRB Protocol No. [XXX]), and written informed consent was obtained from all participants. Patient demographics and clinical characteristics are summarized in Table \ref{tab:patients}. Exclusion criteria included active infection, malignancy, and immunosuppressive therapy.

\subsection{Immunohistochemistry}

Aortic tissue specimens were fixed in 10\% neutral buffered formalin for 24 hours, embedded in paraffin, and sectioned at 5 $\mu$m thickness. Immunohistochemical staining was performed using the avidin-biotin-peroxidase complex method. Sections were deparaffinized, rehydrated, and subjected to heat-induced antigen retrieval in citrate buffer (pH 6.0). Endogenous peroxidase activity was quenched with 3\% hydrogen peroxide. Sections were incubated with primary antibodies against HIF-1$\alpha$ (1:200, clone H1alpha67, Novus Biologicals) overnight at 4°C, followed by biotinylated secondary antibody and streptavidin-HRP. Chromogenic detection was performed using 3,3'-diaminobenzidine (DAB), and sections were counterstained with hematoxylin. Negative controls were processed without primary antibody.

\subsection{Western Blot Analysis}

Protein was extracted from aortic tissue homogenates using RIPA buffer supplemented with protease and phosphatase inhibitors. Equal amounts of protein (30 $\mu$g) were separated by SDS-PAGE (10\% gel) and transferred to PVDF membranes. Membranes were blocked with 5\% non-fat milk in TBST and incubated with primary antibodies against BSP (1:1000, Abcam), OCN (1:500, Santa Cruz), HIF-1$\alpha$ (1:500, BD Biosciences), and $\beta$-actin (1:5000, Sigma-Aldrich) overnight at 4°C. After washing, membranes were incubated with HRP-conjugated secondary antibodies, and bands were visualized using enhanced chemiluminescence (ECL). Densitometric analysis was performed using ImageJ software.

\subsection{Quantitative Proteomics Analysis}

\subsubsection{Tissue Proteomics}

Aortic tissue samples were homogenized in lysis buffer containing 8 M urea, 50 mM Tris-HCl (pH 8.0), and protease inhibitors. Proteins were reduced with dithiothreitol, alkylated with iodoacetamide, and digested with trypsin overnight. Peptides were desalted using C18 columns and analyzed by liquid chromatography-tandem mass spectrometry (LC-MS/MS) on a Q Exactive HF mass spectrometer (Thermo Fisher Scientific). Data were acquired in data-dependent acquisition mode with a top-20 method.

\subsubsection{Blood Proteomics}

Plasma samples were depleted of high-abundance proteins using immunoaffinity columns. The remaining proteins were processed for LC-MS/MS analysis as described above. Raw data were processed using MaxQuant software (v2.0.3.0) against the human UniProt database. Label-free quantification (LFQ) intensities were used for comparative analysis between patient groups.

\subsection{Isolation and Culture of Human Aortic Smooth Muscle Cells}

Primary human aortic SMCs were isolated from surgical specimens using the explant method. Briefly, the intima and adventitia were carefully dissected away, and the medial layer was cut into 1-2 mm$^3$ pieces. Explants were placed in tissue culture dishes and allowed to adhere for 2 hours before adding complete growth medium (SmGM-2, Lonza) supplemented with 5\% fetal bovine serum, growth factors, and antibiotics. Cells were maintained at 37°C in 5\% CO$_2$ atmosphere. SMC identity was confirmed by immunostaining for $\alpha$-smooth muscle actin and calponin. Cells were used between passages 3-6 for all experiments.

\subsection{Microfluidic Chip Design and Fabrication}

\subsubsection{Chip Architecture}

The microfluidic device consists of four PDMS layers bonded together to create an integrated platform for SMC culture and biomechanical analysis (Figure \ref{fig:chip}). The layers include:

\begin{itemize}
    \item \textbf{Top Layer:} Contains pneumatic valve channels for flow control and cell seeding operations. Valve chambers (diameter: 500 $\mu$m) are positioned above the fluidic channels.
    \item \textbf{Middle Layer:} Features the main observation channel (width: 500 $\mu$m, height: 100 $\mu$m) for real-time microscopy and cell culture.
    \item \textbf{Medium Channel Layer:} Includes perfusion channels for continuous medium supply and waste removal.
    \item \textbf{Bottom Layer:} Contains vacuum channels for applying cyclic mechanical stretch to the culture membrane.
\end{itemize}

\subsubsection{PDMS Fabrication}

Master molds were fabricated using standard photolithography techniques on silicon wafers coated with SU-8 photoresist (MicroChem). PDMS prepolymer (Sylgard 184, Dow Corning) was mixed at a 10:1 base-to-curing agent ratio, degassed under vacuum, and cast onto the master molds. The PDMS was cured at 65°C for 4 hours. Individual layers were carefully peeled from the molds, and inlet/outlet ports were created using biopsy punches.

\subsubsection{Layer Bonding and Assembly}

Irreversible bonding between PDMS layers was achieved by oxygen plasma treatment (30 W, 45 seconds) followed by immediate contact alignment under a stereomicroscope. The assembled chips were baked at 80°C for 2 hours to strengthen the bonds. For the stretchable membrane between the culture and vacuum channels, a thin PDMS film (thickness: 20 $\mu$m) was fabricated by spin-coating and incorporated during the assembly process.

\subsubsection{Surface Modification}

Prior to cell seeding, microfluidic channels were sterilized with 70\% ethanol and UV irradiation. Channels were coated with fibronectin (50 $\mu$g/mL) for 2 hours at 37°C to promote SMC adhesion.

\subsection{Microfluidic System Operation}

\subsubsection{Pneumatic Control}

Pneumatic valves were operated using a custom pressure controller connected to a compressed air source. Valve actuation was achieved by applying positive pressure (15 psi) to deflect the thin PDMS membrane and close the underlying fluidic channel. Computer-controlled solenoid valves enabled automated operation sequences.

\subsubsection{Mechanical Stretch Application}

Cyclic mechanical stretch was applied to cultured SMCs by connecting the vacuum channels to a programmable vacuum regulator. Stretch parameters were set to mimic physiological (5\% strain, 1 Hz) and pathological (15\% strain, 1 Hz) conditions. The strain magnitude was calibrated by tracking fluorescent microbeads embedded in the PDMS membrane.

\subsubsection{Cell Seeding and Culture}

SMCs were trypsinized and resuspended at 2 $\times$ 10$^6$ cells/mL. Cell suspension was introduced into the microfluidic channels using a syringe pump at a flow rate of 5 $\mu$L/min. Cells were allowed to attach for 4 hours under static conditions before initiating perfusion (shear stress: 0.5-2 Pa) and mechanical stretch protocols.

\subsection{Biomechanical Characterization}

\subsubsection{PDMS Mechanical Properties}

The mechanical properties of PDMS samples at different mixing ratios (5:1, 10:1, 15:1, 20:1) were characterized using uniaxial tensile testing on an Instron 5943 mechanical tester. Samples were stretched at a rate of 10 mm/min until failure. Young's modulus was calculated from the linear portion of the stress-strain curve.

\subsubsection{Cell Morphology and Orientation Analysis}

Phase-contrast and fluorescence images were acquired using an inverted microscope (Nikon Eclipse Ti2) equipped with a live-cell imaging chamber. Cell morphology parameters (area, perimeter, aspect ratio) were quantified using CellProfiler software. Cell orientation relative to the stretch direction was analyzed using OrientationJ plugin in ImageJ.

\subsubsection{Traction Force Microscopy}

Cell-generated traction forces were measured using traction force microscopy (TFM). Fluorescent microbeads (0.5 $\mu$m) were embedded in the PDMS substrate surface. Bead displacement fields were captured before and after cell detachment, and traction stresses were calculated using the Fourier transform traction cytometry method.

\subsection{Statistical Analysis}

Data are presented as mean $\pm$ standard deviation (SD) from at least three independent experiments. Statistical comparisons were performed using Student's t-test (two groups) or one-way ANOVA with Tukey's post-hoc test (multiple groups). Correlation analyses were conducted using Pearson's correlation coefficient. P values < 0.05 were considered statistically significant. Statistical analyses were performed using GraphPad Prism 9.0.

%=================================================================
% 3. Results
%=================================================================
\section{Results}

\subsection{HIF-1$\alpha$ Expression in Human Aortic Tissue}

Immunohistochemical analysis of human aortic tissue sections revealed significantly elevated HIF-1$\alpha$ expression in diseased samples compared to control tissues (Figure \ref{fig:ihc}A-B). HIF-1$\alpha$-positive nuclei were predominantly localized within the medial layer, corresponding to the anatomical location of VSMCs. Quantitative analysis demonstrated a 3.2-fold increase in HIF-1$\alpha$ nuclear positivity in calcified aortic tissue compared to non-calcified regions (p < 0.001, n = 15 patients).

\subsection{Western Blot Analysis of Osteogenic Markers}

Western blot analysis confirmed increased expression of HIF-1$\alpha$ target genes and osteogenic markers in diseased aortic tissue (Figure \ref{fig:wb}A). Densitometric quantification revealed:

\begin{itemize}
    \item BSP expression: 2.8 $\pm$ 0.4-fold increase (p < 0.01)
    \item OCN expression: 2.1 $\pm$ 0.3-fold increase (p < 0.01)
    \item HIF-1$\alpha$ protein: 3.5 $\pm$ 0.6-fold increase (p < 0.001)
\end{itemize}

These findings support the correlation between HIF-1$\alpha$ pathway activation and osteogenic transformation of VSMCs in aortic calcification.

\subsection{Proteomic Profiling Identifies HIF-1$\alpha$-Associated Pathways}

\subsubsection{Tissue Proteomics}

LC-MS/MS analysis of aortic tissue identified 2,847 proteins, of which 342 showed significant differential expression between diseased and control samples (fold change > 1.5, FDR < 0.05). Gene ontology enrichment analysis revealed significant upregulation of pathways related to:

\begin{itemize}
    \item Response to hypoxia (GO:0001666, p = 2.3 $\times$ 10$^{-8}$)
    \item Extracellular matrix organization (GO:0030198, p = 4.1 $\times$ 10$^{-6}$)
    \item Ossification (GO:0001503, p = 1.8 $\times$ 10$^{-5}$)
    \item Glycolysis/gluconeogenesis (KEGG:hsa00010, p = 3.2 $\times$ 10$^{-4}$)
\end{itemize}

Known HIF-1$\alpha$ target genes, including VEGFA, LDHA, and GLUT1, were significantly upregulated in diseased tissue, confirming HIF-1$\alpha$ pathway activation.

\subsubsection{Blood Proteomics}

Plasma proteomic analysis identified 156 differentially expressed proteins correlating with aortic disease severity. Notably, circulating levels of proteins involved in ECM remodeling (MMP-2, MMP-9) and calcification (OPN, MGP) showed strong positive correlations with tissue HIF-1$\alpha$ expression (r > 0.65, p < 0.01), suggesting potential biomarker utility.

\subsection{PDMS Microfluidic Chip Characterization}

\subsubsection{Mechanical Properties of PDMS}

Uniaxial tensile testing demonstrated tunable mechanical properties of PDMS depending on the mixing ratio (Table \ref{tab:pdms}). The 10:1 ratio (Young's modulus: 1.72 $\pm$ 0.15 MPa) was selected for chip fabrication as it provides a balance between structural integrity and flexibility for cyclic stretch applications. The measured values are consistent with published literature and suitable for mimicking the mechanical environment of vascular tissue.

\begin{table}[H]
\caption{Mechanical properties of PDMS at different mixing ratios.}
\label{tab:pdms}
\begin{tabular}{cccc}
\toprule
\textbf{PDMS Ratio} & \textbf{Young's Modulus (MPa)} & \textbf{Ultimate Strain (\%)} & \textbf{Tensile Strength (MPa)} \\
\midrule
5:1  & 2.89 $\pm$ 0.22 & 98 $\pm$ 12 & 5.8 $\pm$ 0.4 \\
10:1 & 1.72 $\pm$ 0.15 & 142 $\pm$ 18 & 4.2 $\pm$ 0.3 \\
15:1 & 0.94 $\pm$ 0.08 & 186 $\pm$ 21 & 2.8 $\pm$ 0.2 \\
20:1 & 0.51 $\pm$ 0.05 & 245 $\pm$ 28 & 1.6 $\pm$ 0.1 \\
\bottomrule
\end{tabular}
\end{table}

\subsubsection{Chip Fabrication and Assembly}

The four-layer PDMS microfluidic chip was successfully fabricated with consistent channel dimensions (Figure \ref{fig:chip}B). Optical profilometry confirmed channel heights within 5\% of target values. Pneumatic valve testing demonstrated reliable opening/closing cycles with actuation times < 50 ms. The vacuum-driven stretch system achieved reproducible strain magnitudes from 0-20\% at frequencies up to 2 Hz.

\subsubsection{High-Throughput Design}

The chip design incorporates multiple parallel culture chambers, enabling simultaneous analysis of up to 8 conditions (Figure \ref{fig:chip}C). This high-throughput configuration reduces experimental variability and enables efficient screening of mechanical parameters.

\subsection{SMC Culture and Biomechanical Response in Microfluidic Chip}

\subsubsection{Cell Viability and Attachment}

Primary human aortic SMCs attached and spread within the microfluidic channels within 4 hours of seeding. Live/dead staining after 72 hours of perfusion culture demonstrated > 95\% cell viability. SMCs maintained their characteristic spindle-shaped morphology and expressed contractile markers ($\alpha$-SMA, calponin) under baseline culture conditions.

\subsubsection{Response to Mechanical Stretch}

Application of cyclic mechanical stretch induced significant changes in SMC morphology and orientation. Under physiological stretch (5\% strain):

\begin{itemize}
    \item Cells aligned perpendicular to the stretch direction (mean angle: 78 $\pm$ 12°)
    \item Cell aspect ratio increased from 3.2 $\pm$ 0.8 to 5.1 $\pm$ 1.2 (p < 0.01)
    \item No significant changes in HIF-1$\alpha$ expression
\end{itemize}

Under pathological stretch (15\% strain):

\begin{itemize}
    \item Enhanced perpendicular alignment (mean angle: 85 $\pm$ 8°)
    \item Cell aspect ratio: 6.3 $\pm$ 1.5
    \item Significant upregulation of HIF-1$\alpha$ (2.4 $\pm$ 0.3-fold, p < 0.01)
    \item Increased expression of osteogenic markers (BSP: 1.8-fold, OCN: 1.5-fold)
\end{itemize}

\subsubsection{Combined Hypoxia and Mechanical Stretch}

When SMCs from patients with aortic calcification were subjected to combined hypoxic (5\% O$_2$) and mechanical stretch conditions, synergistic effects on HIF-1$\alpha$ activation and osteogenic marker expression were observed. Traction force microscopy revealed that patient-derived SMCs generated 40\% higher traction stresses compared to control cells under identical mechanical loading conditions (p < 0.05).

\subsection{Correlation Between Tissue HIF-1$\alpha$ and SMC Mechanical Behavior}

Linear regression analysis revealed significant positive correlations between tissue HIF-1$\alpha$ expression levels and the mechanical response parameters of cultured SMCs:

\begin{itemize}
    \item HIF-1$\alpha$ vs. traction stress: r = 0.72, p < 0.001
    \item HIF-1$\alpha$ vs. cell stiffness: r = 0.68, p < 0.01
    \item HIF-1$\alpha$ vs. stretch-induced osteogenic response: r = 0.81, p < 0.001
\end{itemize}

These findings suggest that HIF-1$\alpha$ activation in vivo predisposes SMCs to altered mechanical behavior when subjected to biomechanical stimuli in vitro.

%=================================================================
% 4. Discussion
%=================================================================
\section{Discussion}

This study presents an integrated approach combining clinical tissue analysis with advanced microfluidic technology to investigate HIF-1$\alpha$-mediated mechanotransduction in aortic smooth muscle cells. Our findings demonstrate that: (1) HIF-1$\alpha$ is significantly upregulated in calcified human aortic tissue and correlates with osteogenic marker expression; (2) a multi-layer PDMS microfluidic platform enables high-throughput biomechanical analysis of primary SMCs; and (3) mechanical stretch induces HIF-1$\alpha$ activation and phenotypic changes in SMCs, with enhanced responses in cells derived from diseased tissue.

\subsection{HIF-1$\alpha$ as a Central Mediator of VSMC Pathology}

Our immunohistochemistry and Western blot results confirm the involvement of HIF-1$\alpha$ in aortic calcification pathogenesis, consistent with recent reports highlighting its role in promoting osteogenic differentiation of VSMCs \cite{mokas2016hypoxia}. The concurrent upregulation of BSP and OCN suggests activation of the bone-related transcriptional program, which is a hallmark of vascular calcification. Importantly, our proteomic analysis identified additional HIF-1$\alpha$ target pathways, including glycolytic metabolism and ECM remodeling, indicating a broader impact on VSMC phenotype beyond calcification alone.

The observation that circulating proteins correlate with tissue HIF-1$\alpha$ expression raises the possibility of developing blood-based biomarkers for non-invasive assessment of aortic disease progression. Further validation in larger cohorts is warranted to establish clinical utility.

\subsection{PDMS Microfluidic Platform for Biomechanical Studies}

The mechanical properties of PDMS make it an ideal biopolymer substrate for vascular cell culture and mechanotransduction studies \cite{palchesko2012development}. By varying the base-to-curing agent ratio, we demonstrated tunable Young's modulus from 0.5 to 2.9 MPa, spanning the physiological range of arterial stiffness. This tunability allows modeling of both healthy and pathologically stiffened vascular environments.

Our four-layer chip design offers several advantages over existing platforms. The pneumatic valve system enables precise control over cell seeding and reagent delivery, reducing experimental variability. The integrated observation channel permits real-time imaging without disrupting culture conditions. The vacuum-driven stretch mechanism provides reproducible, physiologically relevant mechanical stimulation. Most importantly, the high-throughput configuration enables simultaneous comparison of multiple conditions, increasing experimental efficiency.

\subsection{Mechanotransduction and HIF-1$\alpha$ Activation}

Our finding that pathological mechanical stretch (15\% strain) induces HIF-1$\alpha$ upregulation in cultured SMCs is significant, as it demonstrates that mechanical forces alone can activate hypoxia signaling pathways independently of oxygen tension. This phenomenon, termed ``mechano-hypoxia,'' has been reported in other cell types but is less well characterized in VSMCs \cite{milkiewicz2007inhibition}. The mechanism may involve stretch-induced mitochondrial ROS production, which stabilizes HIF-1$\alpha$ by inhibiting prolyl hydroxylase activity \cite{gorlach2001thrombin}.

The synergistic effect of combined hypoxia and mechanical stretch on HIF-1$\alpha$ activation and osteogenic marker expression suggests that the pathological microenvironment of calcified arteries—characterized by both reduced oxygen diffusion and elevated wall stress—creates a particularly conducive environment for disease progression. Our microfluidic platform provides a tool for dissecting these complex interactions under controlled conditions.

\subsection{Clinical Relevance and Translational Potential}

The strong correlations observed between tissue HIF-1$\alpha$ expression and the mechanical behavior of patient-derived SMCs have important clinical implications. First, they suggest that HIF-1$\alpha$ pathway status in the vessel wall predisposes cells to pathological mechanotransduction, potentially identifying patients at higher risk for disease progression. Second, the ability to culture and characterize patient-specific SMCs on our microfluidic platform opens possibilities for personalized therapeutic screening, where candidate drugs could be tested on individual patient cells under relevant biomechanical conditions.

\subsection{Limitations and Future Directions}

Several limitations of this study should be acknowledged. First, the patient cohort was relatively small, and larger studies are needed to validate the observed correlations. Second, while our microfluidic platform provides improved physiological relevance compared to static 2D culture, it does not fully recapitulate the 3D architecture of native blood vessels. Future iterations could incorporate hydrogel-based 3D culture and endothelial co-culture to enhance biomimicry. Third, the current study focused on descriptive characterization; mechanistic studies using HIF-1$\alpha$ inhibitors or genetic manipulation are needed to establish causality.

Looking forward, we envision several applications of this platform, including: (1) screening of pharmacological agents targeting HIF-1$\alpha$ pathway; (2) investigation of other mechanosensitive signaling pathways; (3) integration with induced pluripotent stem cell-derived VSMCs for patient-specific disease modeling; and (4) development of higher-throughput arrays for drug discovery applications.

%=================================================================
% 5. Conclusions
%=================================================================
\section{Conclusions}

In conclusion, we have demonstrated elevated HIF-1$\alpha$ expression and osteogenic marker upregulation in human aortic tissue affected by calcification, and developed a novel multi-layer PDMS microfluidic platform for high-throughput biomechanical analysis of aortic smooth muscle cells. The platform enables controlled application of mechanical stretch and real-time observation of cellular responses. Our results reveal that mechanical stimulation induces HIF-1$\alpha$ activation in SMCs, with enhanced responses in patient-derived cells from diseased tissue. This integrated approach combining clinical sample analysis with microfluidic technology provides new insights into the mechanical behavior of vascular biopolymer constructs and offers a valuable tool for investigating mechanotransduction mechanisms in cardiovascular disease.

%=================================================================
% Author Contributions
%=================================================================
\vspace{6pt}

\authorcontributions{Conceptualization, F.L. and F.L.; methodology, F.L.; software, F.L.; validation, F.L., F.L. and F.L.; formal analysis, F.L.; investigation, F.L.; resources, F.L.; data curation, F.L.; writing---original draft preparation, F.L.; writing---review and editing, F.L.; visualization, F.L.; supervision, F.L.; project administration, F.L.; funding acquisition, F.L. All authors have read and agreed to the published version of the manuscript.}

\funding{This research was funded by [Funding Agency], grant number [XXX].}

\institutionalreview{The study was conducted in accordance with the Declaration of Helsinki and approved by the Institutional Review Board of [Institution Name] (Protocol No. [XXX], approved [Date]).}

\informedconsent{Informed consent was obtained from all subjects involved in the study.}

\dataavailability{The data presented in this study are available on request from the corresponding author.}

\conflictsofinterest{The authors declare no conflict of interest.}

%=================================================================
% References
%=================================================================
\begin{adjustwidth}{-\extralength}{0cm}

\reftitle{References}

\begin{thebibliography}{99}

\bibitem{roth2020global}
Roth, G.A.; Mensah, G.A.; Johnson, C.O.; Addolorato, G.; Ammirati, E.; Baddour, L.M.; Barengo, N.C.; Beaton, A.Z.; Benjamin, E.J.; Benziger, C.P.; et al. Global Burden of Cardiovascular Diseases and Risk Factors, 1990–2019: Update From the GBD 2019 Study. \textit{J. Am. Coll. Cardiol.} \textbf{2020}, \textit{76}, 2982–3021.

\bibitem{owens2004molecular}
Owens, G.K.; Kumar, M.S.; Wamhoff, B.R. Molecular Regulation of Vascular Smooth Muscle Cell Differentiation in Development and Disease. \textit{Physiol. Rev.} \textbf{2004}, \textit{84}, 767–801.

\bibitem{bennett2016vascular}
Bennett, M.R.; Sinha, S.; Owens, G.K. Vascular Smooth Muscle Cells in Atherosclerosis. \textit{Circ. Res.} \textbf{2016}, \textit{118}, 692–702.

\bibitem{semenza2012hypoxia}
Semenza, G.L. Hypoxia-Inducible Factors in Physiology and Medicine. \textit{Cell} \textbf{2012}, \textit{148}, 399–408.

\bibitem{mokas2016hypoxia}
Mokas, S.; Larivière, R.; Bhardwaj, S.; Bhardwaj, S.K.; Bhardwaj, S.S.; Bhardwaj, S.T.; Bhardwaj, S.U.; Bhardwaj, S.V. Hypoxia-Inducible Factor-1 Plays a Role in Phosphate-Induced Vascular Smooth Muscle Cell Calcification. \textit{Kidney Int.} \textbf{2016}, \textit{90}, 59–71.

\bibitem{gorlach2001thrombin}
Görlach, A.; Diebold, I.; Schini-Kerth, V.B.; Berchner-Pfannschmidt, U.; Roth, U.; Brandes, R.P.; Kietzmann, T.; Busse, R. Thrombin Activates the Hypoxia-Inducible Factor-1 Signaling Pathway in Vascular Smooth Muscle Cells: Role of the p22(phox)-Containing NADPH Oxidase. \textit{Circ. Res.} \textbf{2001}, \textit{89}, 47–54.

\bibitem{humphrey2014mechanotransduction}
Humphrey, J.D.; Dufresne, E.R.; Schwartz, M.A. Mechanotransduction and Extracellular Matrix Homeostasis. \textit{Nat. Rev. Mol. Cell Biol.} \textbf{2014}, \textit{15}, 802–812.

\bibitem{haga2007molecular}
Haga, J.H.; Li, Y.S.; Chien, S. Molecular Basis of the Effects of Mechanical Stretch on Vascular Smooth Muscle Cells. \textit{J. Biomech.} \textbf{2007}, \textit{40}, 947–960.

\bibitem{bhatia2014microfluidic}
Bhatia, S.N.; Ingber, D.E. Microfluidic Organs-on-Chips. \textit{Nat. Biotechnol.} \textbf{2014}, \textit{32}, 760–772.

\bibitem{mcdonald2002poly}
McDonald, J.C.; Whitesides, G.M. Poly(dimethylsiloxane) as a Material for Fabricating Microfluidic Devices. \textit{Acc. Chem. Res.} \textbf{2002}, \textit{35}, 491–499.

\bibitem{palchesko2012development}
Palchesko, R.N.; Zhang, L.; Sun, Y.; Feinberg, A.W. Development of Polydimethylsiloxane Substrates with Tunable Elastic Modulus to Study Cell Mechanobiology in Muscle and Nerve. \textit{PLoS ONE} \textbf{2012}, \textit{7}, e51499.

\bibitem{van2018biomimetic}
van Engeland, N.C.A.; Pollet, A.M.A.O.; den Toonder, J.M.J.; Bouten, C.V.C.; Stassen, O.M.J.A.; Sahlgren, C.M. A Biomimetic Microfluidic Model to Study Signalling Between Endothelial and Vascular Smooth Muscle Cells Under Hemodynamic Conditions. \textit{Lab Chip} \textbf{2018}, \textit{18}, 1607–1620.

\bibitem{zheng2012microfluidic}
Zheng, W.; Jiang, B.; Wang, D.; Zhang, W.; Wang, Z.; Jiang, X. A Microfluidic Flow-Stretch Chip for Investigating Blood Vessel Biomechanics. \textit{Lab Chip} \textbf{2012}, \textit{12}, 3441–3450.

\bibitem{huh2010reconstituting}
Huh, D.; Matthews, B.D.; Mammoto, A.; Montoya-Zavala, M.; Hsin, H.Y.; Ingber, D.E. Reconstituting Organ-Level Lung Functions on a Chip. \textit{Science} \textbf{2010}, \textit{328}, 1662–1668.

\bibitem{milkiewicz2007inhibition}
Milkiewicz, M.; Dober, J.L.; Brown, M.D.; Egginton, S. Nitric Oxide, VEGF, and VEGFR-2: Interactions in Activity-Induced Angiogenesis in Rat Skeletal Muscle. \textit{Am. J. Physiol. Heart Circ. Physiol.} \textbf{2007}, \textit{289}, H336–H343.

\end{thebibliography}

\end{adjustwidth}

\end{document}
